\ruleblock{Quantificatore Universale Dipendente ($\forall$)}{
    \begin{minipage}{0.48\textwidth}
        \centering
        \begin{prooftree}
            \hypo{x:\sigma \vdash \Phi} 
            \infer1[($\forall I$)]{\forall x^\sigma \Phi}
        \end{prooftree}
    \end{minipage}%
    \hfill
    \begin{minipage}{0.48\textwidth}
        \centering
        \begin{prooftree}
            \hypo{\forall x^\sigma \Phi} 
            \hypo{t:\sigma} 
            \infer2[($\forall E$)]{\Phi[t/x]}
        \end{prooftree}
    \end{minipage}
}{Nel caso dipendente, la regola di introduzione richiede che il termine $x$ sia 
di tipo $\sigma$. La regola di eliminazione utilizza un testimone $t:\sigma$ per 
istanziare il predicato $\Phi$, risultando nella sostituzione $\Phi[t/x]$.}

\ruleblock{Quantificatore Esistenziale Dipendente ($\exists$)}{
    \begin{minipage}{0.42\textwidth}
        \centering
        \begin{prooftree}
            \hypo{t:\sigma} 
            \hypo{\Phi[t/x]} 
            \infer2[($\exists I$)]{\exists x^\sigma \Phi}
        \end{prooftree}
    \end{minipage}%
    \hfill
    \begin{minipage}{0.55\textwidth}
        \centering
        \begin{prooftree}
            \hypo{\exists x^\sigma \Phi} 
            \hypo{x:\sigma, [\Phi]^i \vdash \Psi} 
            \infer2[($\exists E^i$)]{\Psi}
        \end{prooftree}
    \end{minipage}
}{L'introduzione esistenziale richiede sia il testimone $t$ che il suo tipo $\sigma$. 
Per l'eliminazione, la variabile $x$ di tipo $\sigma$ viene scaricata insieme all'ipotesi 
$\Phi$ per raggiungere la conclusione $\Psi$.}